\begin{table}[H]
\centering
\caption*{\textbf{Quadro B.1 - Requisitos de conformidade da MOSAIC 0.2}}
\scriptsize
\renewcommand{\arraystretch}{1.10}
\begin{tabularx}{\linewidth}{|p{0.07\linewidth}|Y|p{0.10\linewidth}|}
\hline
\textbf{ID} & \textbf{Requisito} & \textbf{Obrigatório} \\
\hline
C01 & O planner inicial recebe o pedido sem nomes de skills. & Sim \\
\hline
C02 & Cada nó possui resultado observável, \texttt{doneWhen} e participa de um DAG válido. & Sim \\
\hline
C03 & O feedback do catálogo ocorre somente após P0 e no máximo uma vez antes da execução. & Sim \\
\hline
C04 & O reranking lê o body canônico das candidatas. & Sim \\
\hline
C05 & O bundle admite cardinalidade de zero a \(K_{\max}\) e ordem determinística. & Sim \\
\hline
C06 & O menu é \(T_{\mathrm{base}}\cup\bigcup A(s)\) para as skills selecionadas. & Sim \\
\hline
C07 & A decisão do modelo contém somente campos semânticos e avaliações ordenadas dos critérios. & Sim \\
\hline
C08 & O runtime captura automaticamente cada retorno bem-sucedido de tool executável na ordem dos resultados. & Sim \\
\hline
C09 & Dados de chamada e retorno ausentes, duplicados ou não correlacionados causam falha de runtime. & Sim \\
\hline
C10 & A tool terminal de saída estruturada nunca é registrada como observação. & Sim \\
\hline
C11 & \texttt{needs\_revision} exige ao menos uma observação e associa o pedido ao conjunto completo do nó. & Sim \\
\hline
C12 & O planner localizado recebe nome, entrada e saída de todas as observações, sem \texttt{callId}. & Sim \\
\hline
C13 & Revisões preservam nós concluídos e seus resultados. & Sim \\
\hline
C14 & A resposta final é empacotada sem nova chamada cognitiva. & Sim \\
\hline
C15 & Nenhum objeto estruturado escrito pelo modelo altera plano ou estado antes de satisfazer o contrato completo do runtime. & Sim \\
\hline
\end{tabularx}
\par\smallskip\raggedright\small Fonte: elaboração própria (2026).
\end{table}
